\section{Problem and motivation}

Policy analysis of natural gas markets often rests on large equilibrium models, and these models need substantial expert work to stay in use.
The Global Gas Model (GGM) \citep{egging2019global} is one of them.
It descends from earlier equilibrium models of gas markets \citep{egging2010world} and is used for questions such as how much Russian gas the European Union still needs under different climate pathways \citep{barner2025russian}.
GGM represents the natural gas value chain from production, through pipelines or LNG, to consumption and storage, for more than ninety countries, and lets suppliers behave anywhere between perfect competition and Cournot competition.
The data for a scenario are held in three spreadsheet workbooks: one describing the world (infrastructure, costs, market power), one with projected production and consumption, and one with calibrated values.

Such a model is only as credible as the work done around it.
Its input data are collected from statistical sources and published outlooks, a scenario is defined as changes to those inputs, and the model is calibrated so that it reproduces observed and projected market outcomes.
Calibrating it for a new scenario takes an experienced analyst at least several days, and up to several weeks after a major data revision, because the analyst ``needs to combine market knowledge with model expertise and trade off choices'' \citep[p.~26]{egging2019global}.
Much of how this is done is not written down and is held by a few people (Chapter 3).

Agents built on large language models (LLMs) differ from earlier automation in that they can work with unstructured material: they read reports and documentation, write and run code, and call other software.
They are now used to set up, calibrate and operate existing models in fields from nuclear engineering to energy markets \citep{abulawi2026autosam, gomezguillen2026agentic, golden2026agentic}, with uneven results that Chapter 2 reviews.
Much of the work around GGM is of the kind they have been built for: finding figures in reports, mapping them to the model's countries and regions, and applying documented rules.

How this transfers to GGM is not obvious.
Some of the work is mechanical and better done by ordinary code, whose rules give one correct answer.
Some rests on judgment, and calibration is the clearest case: its parameters, such as market power, are not measured anywhere and are not pinned down by the data, so setting them is partly an economic judgment.
Where judgment is involved, the risk is not a run that fails, but a plausible result resting on a poor choice, which an agent's fluent explanation could make more convincing rather than less.
Which work should go to an agent, which to code and which to a person is therefore the question this thesis addresses.
